% The actual content of the blueprint, used by both the web and print versions.
% Do not include \begin{document}.
%
% Each node carries:
%   \lean{Full.Declaration.Name}  the Lean declaration this node IS
%   \leanok                        this node is fully formalized
%   \uses{label, ...}              its immediate dependencies
%
% `leanblueprint checkdecls` verifies every \lean{...} names a declaration that actually exists
% in the built project. That check is what keeps this document from drifting away from the Lean,
% so never write a \lean{} you have not built, and run it before every push.

\chapter{Frankl-Kupavskii-Swanepoel Theorem 3: fewer than C(d+2,2) edges embeds in R^d}

\begin{definition}[Pairwise coprimality]
  \label{def:pairwise-coprime}
  \lean{FKSEdgeThreshold.Standalone.Mathlib.PairwiseCoprime}
  \leanok
  A finite set of naturals is \emph{pairwise coprime} when any two distinct elements are coprime.
\end{definition}

\begin{lemma}[Small primes are pairwise coprime]
  \label{lem:small-primes-coprime}
  \lean{FKSEdgeThreshold.Standalone.Mathlib.SmallPrimesCoprime}
  \leanok
  \uses{def:pairwise-coprime}
  The set $\{2, 3, 5\}$ is pairwise coprime.
\end{lemma}

\begin{proof}
  \leanok
  \uses{def:pairwise-coprime}
  Each pair is distinct and shares no prime factor, so the claim is decidable and holds by
  computation.
\end{proof}
